\documentclass[11pt,a4paper]{article}
\usepackage[margin=2.1cm]{geometry}
\usepackage[T1]{fontenc}
\usepackage{lmodern}
\usepackage{xcolor}
\usepackage{booktabs}
\usepackage{tabularx}
\usepackage{longtable}
\usepackage{enumitem}
\raggedbottom
\usepackage{needspace}
\usepackage[hidelinks,colorlinks=true,linkcolor=febblue,urlcolor=febblue]{hyperref}
\definecolor{febblue}{HTML}{1C4687}
\setlength{\parskip}{4pt}
\setlist{nosep}

\title{\color{febblue}\LARGE RC Autonomous Car: Budget and Options\\[4pt]\large FEB Autonomous Team}
\author{Formula Electric at Berkeley}
\date{September 2026}

\begin{document}
\maketitle

\section*{The ask}

One 1:10-scale autonomous car on the RoboRacer (F1TENTH) standard: \textbf{\$1{,}550 now},
a complete car on the lidar and camera we already own, plus \textbf{\$400 for the lidar
upgrade} when funds allow (\$1{,}950 together); \$3{,}250 for the official reference build.
Every item is listed with its shop link in Section~\ref{sec:bom}. The FSAE car is on a two-year cycle and cannot host autonomous work this
season. The RC car is where the team develops, tests and learns in the meantime, and the
software pipeline built on it is the one that goes onto the FSAE car next year.

\Needspace{14\baselineskip}
\section{What the car is and how it drives}

\begin{tabularx}{\linewidth}{lX}
\toprule
Part & Role \\
\midrule
Chassis & Traxxas Fiesta ST Rally VXL, ready to run: Slash 4x4 chassis, 4WD, brushless motor, steering servo. The chassis the current RoboRacer build guide and deck drawings are made for, and the car our simulator models. \\
Motor controller & VESC-type. Drives motor and servo from USB and reports wheel rotation back as speed. \\
Computer & Runs ROS~2: sensor drivers, perception, planning, control. \\
2D lidar & A range per angle, 5 to 40 times a second. Sees walls and cones; enough to drive with no map. \\
Camera & Colour: which cone is which. Needed for the perception work. \\
Joystick & Bluetooth gamepad for manual driving and the dead-man switch: the car drives autonomously only while a button is held. \\
Power & One 3S lithium pack drives the motor; the computer takes pack voltage through a fuse; a 5~V converter feeds lidar and camera. \\
Deck & Laser-cut plates and standoffs from the published drawings. \\
\bottomrule
\end{tabularx}

\smallskip
\textbf{Flow:} lidar, wheel odometry and camera in; perception locates walls or cones; a
planner picks a bearing and speed; a controller turns them into steering and throttle; the
VESC actuates, gated by the dead-man. Same nodes and topic names as the team's simulator, so
a stack developed there drops onto the car with only the driver nodes swapped.

\Needspace{14\baselineskip}
\section{Budget}

\begin{tabularx}{\linewidth}{Xrrr}
\toprule
Item & Now & With S2 & Reference \\
\midrule
Chassis, batteries, charger, bag & 560 & 560 & 560 \\
Computer: Jetson Orin Nano Super, NVMe, display plug, antenna, cable & 522 & 522 & 522 \\
Motor controller: Flipsky Mini FSESC 4.20 / VESC 6 MkVI, cables, adapters & 112 & 112 & 276 \\
Power: RoboRacer powerboard, pigtail, balance extender & 164 & 164 & 164 \\
Joystick & 30 & 30 & 30 \\
Lidar: RPLidar A1M8 (owned) / RPLidar S2 / Hokuyo UST-10LX & 0 & 399 & 1{,}200 \\
Camera: RealSense D435i (owned) / RealSense D435i & 0 & 0 & 334 \\
Deck (cut on campus), standoffs, fastener kits & 101 & 101 & 101 \\
Spares allowance & 60 & 60 & 60 \\
\midrule
\textbf{Total} & \textbf{1{,}549} & \textbf{1{,}948} & \textbf{3{,}247} \\
\bottomrule
\end{tabularx}

\smallskip
Prices are September 2026 shop prices (Section~\ref{sec:bom}). \textbf{Now} is the
official build with the sensors we own and a cheaper motor controller: a complete car that
drives laps at 1 to 2~m/s. \textbf{With S2} adds the time-of-flight RPLidar, which bolts on
later without any other change; it is the one upgrade that matters. \textbf{Reference} is
the official RoboRacer bill of materials at today's prices (NVIDIA raised Jetson prices by
60 to 100\% in July 2026, which is why the computer costs \$399, not \$249).

\Needspace{14\baselineskip}
\section{Exact bill of materials}
\label{sec:bom}

Every line links to the shop page. Prices are what the pages showed on 24 September 2026
(Amazon prices from the official RoboRacer sheet of May 2026, which we checked against
the listings). ``Now'' is the car we can build today; ``S2'' is the lidar upgrade line.
The official sheet: \href{https://docs.google.com/spreadsheets/d/1WT-UdjLTnyrb7CvHcuCxKLnRIN8DWroQ/edit}{RoboRacer master BOM}.

\small
\begin{longtable}{p{11.6cm}lr}
\toprule
Item & Shop & USD \\
\midrule
\endfirsthead
\toprule
Item & Shop & USD \\
\midrule
\endhead
\multicolumn{3}{l}{\textbf{Chassis and batteries}} \\
\href{https://traxxas.com/products/landing/fiesta-st-rally-vxl/}{Traxxas Ford Fiesta ST Rally VXL, 74276-4} (backordered at some dealers; Amazon B0D5M8D2R8 also lists it) & Traxxas & 429.95 \\
\href{https://zeeebattery.com/products/zeee-3s-lipo-battery-6000mah-11-1v-80c-ec5}{Zeee 3S 6000~mAh 80C hard case, EC5, 2-pack} (\href{https://a.co/d/07viD7OJ}{Amazon} 89.99) & Zeee & 88.99 \\
\href{https://a.co/d/04M3SdSJ}{LiPo balance charger with EC5 lead} & Amazon & 33.99 \\
\href{https://a.co/d/0gqCuJEP}{Zeee fireproof LiPo bag} & Amazon & 6.99 \\
\multicolumn{3}{l}{\textbf{Computer}} \\
\href{https://www.sparkfun.com/nvidia-jetson-orin-nano-developer-kit.html}{NVIDIA Jetson Orin Nano Super Developer Kit} (backorder; \href{https://www.seeedstudio.com/NVIDIAr-Jetson-Orintm-Nano-Developer-Kit-p-5617.html}{Seeed} 450.30; \href{https://www.amazon.com/dp/B0BZJTQ5YP}{Amazon}); Wi-Fi card and antennas included & SparkFun & 399.00 \\
\href{https://a.co/d/0fV5xrP8}{NVMe SSD, 500~GB, M.2 2280} (sheet price; street price about 60) & Amazon & 100.00 \\
\href{https://www.amazon.com/dp/B071CGCTMY}{DisplayPort dummy plug} (headless boot) & Amazon & 6.99 \\
\href{https://www.amazon.com/dp/B07V1W759Z}{7~dBi Wi-Fi antenna with MHF4 cable} (range; the mount is a campus print) & Amazon & 7.99 \\
\href{https://a.co/d/0hjc5X8u}{Short USB-A to micro-USB cable, Jetson to controller} & Amazon & 7.99 \\
\multicolumn{3}{l}{\textbf{Motor controller}} \\
\href{https://www.amazon.com/dp/B08725X8CT}{Flipsky Mini FSESC 4.20 50~A} (3S to 13S; \href{https://flipsky.net}{flipsky.net} 85, on sale 56) & Amazon & 85.00 \\
\href{https://www.adafruit.com/product/3893}{JST-PH 3-pin to header cable} (controller PPM port to steering servo) & Adafruit & 1.25 \\
\href{https://www.amazon.com/dp/B076F64ZCJ}{Header pins} (three are used) & Amazon & 8.99 \\
\href{https://www.amazon.com/dp/B07Z4FRQNX}{Bullet adapters, 4~mm male to 3.5~mm female, 3-pack} (Traxxas motor leads are 3.5~mm; match the controller's bullets on arrival) & Amazon & 7.99 \\
\href{https://www.amazon.com/dp/B09MQ1MJRD}{EC5 female to XT60 male adapters, 3-pack} (battery to controller) & Amazon & 8.99 \\
\multicolumn{3}{l}{\textbf{Power distribution}} \\
\href{https://orders.ambimat.com/product/roboracer-power-board-1-qty-non-rohs-compliant/}{RoboRacer power board, 1 unit} (19~V to the Jetson, 12~V and 5~V rails, switch, fused; 125 each for three or more) & Ambimat & 145.00 \\
\href{https://a.co/d/08r6C0zE}{Barrel jack to pigtail, 2.5~x~5.5~mm} (power board to Jetson) & Amazon & 9.99 \\
\href{https://a.co/d/04z1CrzG}{3S balance plug extender} (battery to power board) & Amazon & 8.99 \\
\multicolumn{3}{l}{\textbf{Manual control}} \\
\href{https://a.co/d/04eI03bz}{8BitDo Ultimate 2C wireless gamepad} (dead-man and manual drive) & Amazon & 29.99 \\
\multicolumn{3}{l}{\textbf{Sensors}} \\
RPLidar A1M8, owned; \href{https://www.slamtec.com/en/Lidar/A1}{spec} & -- & 0 \\
Intel RealSense D435i, owned; \href{https://store.realsenseai.com/buy-intel-realsense-depth-camera-d435i.html}{spec} & -- & 0 \\
\multicolumn{3}{l}{\textbf{Deck and hardware}} \\
\href{https://drive.google.com/drive/folders/1NU4FZzvMEGKCOFzDBvnjyePnnSMvsZPG}{Platform deck, laser cut from the official DXF} (3~mm acrylic sheet, campus cutter) & material & 20.00 \\
\href{https://drive.google.com/drive/folders/1sy-XiJJ4hmhEKf5qQbUaPYY6Aw-L31Gk}{Antenna mount} and A1M8 adapter plate, campus 3D printer & print & 0 \\
\href{https://a.co/d/0ePVll7A}{M3 socket-head screw kit} & Amazon & 22.86 \\
\href{https://www.amazon.com/dp/B07SBQ1JX5}{M5 socket-head screw kit} (two M5~x~20 are used) & Amazon & 8.66 \\
\href{https://www.mcmaster.com/94868A009}{M3 hex standoffs, 14~mm} & McMaster & 3.68 \\
\href{https://www.mcmaster.com/94868A013}{M3 hex standoffs, 25~mm} & McMaster & 10.76 \\
\href{https://www.mcmaster.com/94868A190}{M3 hex standoffs, 45~mm} & McMaster & 27.90 \\
\href{https://www.mcmaster.com/94500A213/}{M2.5~x~6~mm screws} (Jetson) & McMaster & 5.30 \\
\href{https://www.mcmaster.com/95947A005/}{M2.5 standoffs, 10~mm} (Jetson) & McMaster & 2.16 \\
\multicolumn{3}{l}{\textbf{Spares}} \\
Allowance: bumper, servo saver, tyres, pinion, ordered from the \href{https://traxxas.com/products/parts}{Traxxas parts finder} as they break & Traxxas & 60.00 \\
\midrule
\textbf{Now} & & \textbf{1{,}549} \\
\midrule
\multicolumn{3}{l}{\textbf{Lidar upgrade}} \\
\href{https://www.seeedstudio.com/RPLiDAR-S2-Low-Cost-360-Degree-Laser-Range-Scanner-30M-Range-p-5198.html}{RPLidar S2, 30~m time of flight} (354.99 on sale today; \href{https://www.robotshop.com/products/rplidar-s2-360-laser-scanner-30-m}{RobotShop} also stocks it) & Seeed & 398.99 \\
\midrule
\textbf{Now plus S2} & & \textbf{1{,}948} \\
\bottomrule
\end{longtable}
\normalsize

\smallskip
\textbf{Fitting \$1{,}800 with the S2.} At today's sale prices (S2 at 355, controller at 56)
the total is \$1{,}875. Replacing the power board with a
\href{https://www.pololu.com/product/2851}{Pololu 5~V 5~A regulator} (32.95) for the lidar,
with the Jetson fed from the balance extender through the barrel pigtail, saves \$112 and
lands at \$1{,}760 to \$1{,}835; it costs one evening of soldering and loses the fuse and switch.
Cheapest of all is buying the S2 as a second, later request: the car does not wait for it.

\textbf{What the A1M8 costs us.} Same ROS message, same software, a printed plate instead of
the Hokuyo bolt pattern, USB power instead of 12~V, so nothing about the build gets harder.
What changes is the ceiling: 8{,}000 samples a second at 5 to 10~Hz gives about 1\textdegree{}
resolution, ranges reliable to roughly 6~m indoors, and a scan that takes 100 to 180~ms while
the car moves, so the laps are 1 to 2~m/s and the maps are coarse. Cones at 6~m give two to four
returns, enough for the cone pipeline. Triangulation lidars also wash out in direct sunlight,
so the A1M8 car stays indoors. The S2 lifts all of that (30~m, 32{,}000 samples a second, works
outdoors) and is the one part that makes the car fast.


\Needspace{14\baselineskip}
\section{Options in each area}

\Needspace{10\baselineskip}
\subsection*{Lidar}
The lidar sets the speed the car can drive: at 10~Hz it moves 25~cm between scans at 3~m/s.
All publish the same ROS message, so a swap changes no code; the RPLidar sits on a small
printed adapter plate and the Hokuyo bolts straight to the deck.

\begin{tabularx}{\linewidth}{p{3.2cm}rrrX}
\toprule
Lidar & USD & Rate & Range & Reality \\
\midrule
\href{https://www.slamtec.com/en/Lidar/A1}{RPLidar A1M8} (owned) & 0 & 5 to 10~Hz & 12~m & laps at 1 to 2~m/s; 1\textdegree{} resolution, coarse maps; runs the car until the S2 arrives \\
\href{https://www.slamtec.com/en/Lidar/A2}{RPLidar A2M12} & 229 & 12~Hz & 12~m & fine indoors; noisy on dark surfaces \\
\href{https://www.slamtec.com/en/Lidar/S2}{RPLidar S2} & 399 & 10~Hz & 30~m & time of flight; clean maps; to 3~m/s \\
\href{https://www.slamtec.com/en/Lidar/A3}{RPLidar A3} & 599 & 15 to 20~Hz & 25~m & fastest cheap unit \\
Hokuyo UST-10LX & 1{,}200 & 40~Hz & 10~m & the reference; what race winners run \\
\bottomrule
\end{tabularx}

\Needspace{10\baselineskip}
\subsection*{Compute}
\begin{tabularx}{\linewidth}{p{4.6cm}rXX}
\toprule
Computer & USD & On the RC car & On the FSAE car \\
\midrule
\href{https://www.raspberrypi.com/products/raspberry-pi-5/}{Raspberry Pi 5, 8~GB} & 130 & lidar-only driving; no camera perception & no \\
\href{https://www.sparkfun.com/products/22098}{Jetson Orin Nano Super, 8~GB} & 399 + 100 NVMe & the official part; fits the deck; enough for everything the car does & no: a 3D lidar and several cameras need more \\
\href{https://www.nvidia.com/en-us/autonomous-machines/embedded-systems/jetson-orin/}{Jetson Orin NX 16~GB} on the Nano carrier & about 1{,}200 & fits; overkill & a first driverless stack: M1 lidar, two cameras, detector \\
Jetson AGX Xavier 16~GB (owned, 2018) & 0 & no: 630~g box, does not fit the deck & a bridge for one or two seasons; last software is JetPack~5 \\
\href{https://www.nvidia.com/en-us/autonomous-machines/embedded-systems/jetson-orin/}{Jetson AGX Orin 64~GB} & 3{,}500 & no: 1.4~kg & the full answer, when the stack needs it \\
\bottomrule
\end{tabularx}

\smallskip
Recommended: the Orin Nano on the RC car, permanently. The FSAE computer is decided next
year with a season of data; the owned Xavier carries FSAE perception development until then.
If the club prefers one purchase for both cars, the Orin NX is the largest that fits the RC
car and moves to the FSAE car later, at about \$800 more, leaving the RC car needing a Nano then.

\textbf{Why a Jetson from day one.} Jetson to Jetson is the same software: same operating
system image, same ARM64 build, same CUDA and TensorRT, same ROS~2 packages. Moving the stack
from the Orin Nano to an Orin NX or AGX Orin is copying it across. Developing on an x86 laptop
or PC instead means a different architecture, different driver and CUDA versions, and
recompiling everything, with the porting landing at the worst moment, when the FSAE car is
finally ready. Building on the Nano now makes the later upgrade trivial.

\textbf{A separate line for the on-car FSAE computer.} If the budget process allows it, the
FSAE car's computer should be its own item, requested once the team has a season of data on
what the stack needs: about \$1{,}200 for an Orin NX 16~GB on a carrier board or about
\$3{,}500 for an AGX Orin developer kit, at the July 2026 prices. It is not part of this ask,
and buying it now would mean guessing. Two-year total, RC car with the S2 now plus the FSAE
computer next year: about \$3{,}150 with an Orin NX, about \$5{,}450 with an AGX Orin.

\Needspace{10\baselineskip}
\subsection*{Camera}
The \textbf{RealSense D435i we own} goes on the car: colour for cone classes, depth as a
bonus, and it is the camera in the official bill of materials. Nothing to buy. If a second
camera is ever needed: \href{https://www.amazon.com/dp/B006JH8T3S}{Logitech C920}, \$75,
rolling shutter, fine at these speeds;
\href{https://shop.luxonis.com/products/oak-d-lite-1}{OAK-D Lite}, \$269, stereo depth and
on-camera inference.
\href{https://shop.luxonis.com/products/oak-d-pro-w}{OAK-D Pro W} or
\href{https://store.stereolabs.com/products/zed-2i}{ZED 2i}, \$400 to \$500, global shutter
and weatherproof: FSAE-grade cameras that also fit the RC car, if the club wants the camera
bought once.

\Needspace{10\baselineskip}
\subsection*{Controller}
The controller must accept the car's 3S pack (9 to 12.6~V). The
\href{https://www.amazon.com/dp/B08725X8CT}{Flipsky Mini FSESC 4.20 50~A}, \$85 (\$56 on sale
at \href{https://flipsky.net}{flipsky.net}), runs 3S to 13S on VESC~4.12 hardware, the design the
original F1TENTH cars raced on; USB cable included. The
\href{https://trampaboards.com/vesc-6-mkvi--12s-c-1630.html}{VESC 6 MkVI}, \$250, is the
reference part: better cooling and support. Note that the popular Flipsky Mini FSESC 6.7
is rated 14~V minimum (4S) and does \emph{not} suit this car. Either listed unit is a drop-in.

\Needspace{14\baselineskip}
\section{What we already own, and where it goes}

\begin{itemize}
  \item \textbf{RPLidar A1M8.} Onto the RC car: laps at 1 to 2~m/s until the S2 arrives, then a spare. USB-powered, four screws on a printed plate.
  \item \textbf{RoboSense M1 3D lidar.} It can be bolted to the RC car on a custom plate, and the car would drive, badly: 0.7~kg on a 3~kg car spoils the handling, its half-metre blind zone hides the nearest wall in a 2~m corridor, and its 120\textdegree{} forward view gives wall-following nothing to work with. It is an FSAE sensor; with the Xavier it forms a perception rig for cone-track data this semester.
  \item \textbf{AGX Xavier 16~GB.} Computer for that rig and a bridge for the FSAE car; not for the RC car.
  \item \textbf{RealSense D435i.} The RC car's camera (it is the official part); the rig borrows it when the car is idle.
  \item \textbf{Tamiya 1:10 chassis.} Motorless touring kit: no drawings, no match to the simulator, a semester of mechanical work to save \$300. Bench mule.
  \item \textbf{Alameda test track.} Outdoor sessions for the rig and, later, the FSAE car; not needed for the RC car, which runs in a corridor.
\end{itemize}


\Needspace{14\baselineskip}
\section{Why it is worth it}

\Needspace{10\baselineskip}
\subsection*{This semester}
\begin{itemize}
  \item The FSAE car is being built and cannot host autonomy; without a platform the team has nothing to run on until next year.
  \item Members stay for what they can see move.
  \item Real dynamics and real sensors cannot be simulated: grip, servo backlash, battery sag, a lidar that sees chair legs.
  \item New members learn wiring, power, safety, motor control and ROS~2 on hardware at low cost and low risk.
  \item The simulator and the car share one codebase, so every member's stack has somewhere to go.
\end{itemize}

\Needspace{10\baselineskip}
\subsection*{Carrying over to the FSAE car}
\begin{itemize}
  \item \textbf{Same pipeline.} Cone detection with camera plus lidar, midline between cones, raceline, tracking controller: the Formula Student Driverless stack at small scale.
  \item \textbf{Same tools.} ROS~2 nodes and topics, sensor drivers, bag recording and replay, sim-to-real testing, safety procedures.
  \item \textbf{Same method.} Develop in the simulator, transfer to the car, measure the gap, close it.
  \item \textbf{What does not transfer:} vehicle dynamics and tuned parameters. A 3~kg car says nothing about a 250~kg car's tyres.
\end{itemize}

\Needspace{14\baselineskip}
\section{Will it work, and the practical realities}

The build is the F1TENTH/RoboRacer platform, assembled by hundreds of university teams since
2016 from a public build guide; the parts list is that guide's own with the sensors we own
and a cheaper controller. Driving laps is a solved problem, and our simulator already runs the same
car and software. The car can run the whole pipeline: mapping, localisation, a raceline and a
proper controller; the cheaper sensors cap speed, not capability. The recommended car is
eligible for RoboRacer physical events as built; a lidar upgrade is what would make it fast.

\begin{itemize}
  \item \textbf{One owner.} The car needs a named person for assembly, batteries and access, or it goes unused. This is the main risk.
  \item \textbf{Time.} Two to four weekends to assemble and integrate; first autonomous lap within a month; mapping, racelines and perception fill the semester.
  \item \textbf{Speed.} 1 to 2~m/s on the A1M8, about 3~m/s on the S2, faster only with the Hokuyo. The simulator car goes faster; members will feel the gap.
  \item \textbf{Access.} One car; time on it is earned through the simulator's qualification rule.
  \item \textbf{Batteries.} LiPo bag, supervised charging, a written procedure before the first charge.
  \item \textbf{Wear.} Bumpers and servo savers break; the spares line covers it.
  \item \textbf{Not in this budget.} A 1:5-scale car (Traxxas X-Maxx, Arrma Kraton 8S, Losi 5IVE-T: about \$3{,}500 even with the M1 and Alameda counted, and not runnable on campus), a second car, and the FSAE computer. All later decisions.
\end{itemize}

\bigskip
\noindent RoboRacer build guide: \url{https://f1tenth.readthedocs.io}; official master BOM: \href{https://docs.google.com/spreadsheets/d/1WT-UdjLTnyrb7CvHcuCxKLnRIN8DWroQ/edit}{Google Sheet}\\
Team simulator and manuals: \url{https://github.com/Pranman1/feb-racing}

\end{document}
